\textbf{Цели работы:}
\begin{itemize}
	\item познакомиться с технологией Git и её концепциями;
	\item изучить основные команды и возможности git.
\end{itemize}

\textbf{Задачи:}
\begin{enumerate}
	\item изучить методические материалы к лабораторной работе;
	\item установить git;
	\item создать новый репозиторий, добавить файлы и создать коммит;
	\item создать и настроить мастер репозиторий, запушить в него изменения;
	\item склонировать мастер репозиторий;
	\item разрешить конфликт версий файла.
\end{enumerate}

\section*{Выполнение задания}

\begin{figure}[!htbp]
	\centering
	\includegraphics[width=0.6\linewidth]{"Screenshot 2026-04-09 024951"}
	\caption{Git Bash}
\end{figure}
\FloatBarrier

\VerbatimInput[fontsize=\footnotesize, frame=single, numbers=left, numbersep=5pt, breaklines, breakanywhere=true]{assets/output.txt}
Листинг 1 -- Вывод команд

\section*{Заключение}

В ходе выполнения лабораторной работы были решены следующие задачи:
\begin{itemize}
	\item проведено ознакомление с технологией Git и её концепциями;
	\item изучены основные команды и возможности git.
\end{itemize}

\section*{Контрольные вопросы}

\begin{enumerate}
	\item \textit{Что такое Git?} Распределенная система контроля версий с открытым исходным кодом.
	\item \textit{Что такое ветка?} Перемещаемый указатель коммит.
	\item \textit{Объясните метод объединения веток fast-forward.} Указатель переносится вперед, так как ветка main не имеет параллельных версий, а накатываемые изменения линейны.
	\item \textit{Объясните метод объединения веток merge.} Выполняется простое трёхстороннее слияние, используя последние коммиты объединяемых веток и общего для них родительского коммита, и создается новый merge-коммит.
\end{enumerate}
